\section{Related Work}

\subsection{Code Generation Benchmarks}
The rapid evolution of Code Large Language Models (LLMs)—spanning open-source models like CodeLLama~\citep{roziere2023code}, DeepSeek-Coder~\citep{zhu2024deepseek}, and Qwen-Coder~\citep{hui2024qwen2}, as well as proprietary systems like Claude~\citep{claude4} series, GPT~\citep{gpt4o,GPT-4.1} series, and the Gemini~\citep{gemini2.5} family—has fundamentally reshaped software development. This progress necessitates robust and contemporary benchmarks to accurately assess their capabilities in tasks such as code generation and debugging. Pioneering benchmarks like HumanEval~\citep{chen2021evaluatinglargelanguagemodels} and MBPP~\citep{austin2021programsynthesislargelanguage} established a foundation by evaluating functional correctness on short, algorithm-centric Python problems. However, they are hampered by limitations such as potential data contamination, narrow programming language coverage, and a disconnect from real-world applications. To overcome these shortcomings, subsequent benchmarks~\citep{jimenez2024swebench,jain2025livecodebench,wang2025ojbench,guo2025codeeditorbenchevaluatingcodeediting,zhang2025codecriticbenchholisticcodecritique,zhang2025artifactsbench} have targeted more complex programming tasks. A notable line of work focuses on algorithmic challenges from programming contests~\citep{hendrycks2021measuring,Li_2022,jain2025livecodebench,wang2025ojbench,zheng2025livecodebenchproolympiadmedalists}. LiveCodeBench~\citep{jain2025livecodebench} mitigates data contamination by continuously sourcing new problems from competitive programming platforms. OJBench~\citep{wang2025ojbench} presents a rigorous, competition-level benchmark with 232 problems from prestigious contests like NOI and ICPC, demanding advanced code reasoning. APPS~\citep{hendrycks2021measuring} offers a large-scale dataset of 10,000 problems categorized by difficulty. Another stream of research evaluates more comprehensive and multilingual capabilities~\citep{cassano2022multiplescalableextensibleapproach,peng2024humanevalxlmultilingualcodegeneration,zhang2024naturalcodebenchexaminingcodingperformance,chai2025mceval,fullstackbench}. McEval~\citep{chai2025mceval} is a massively multilingual benchmark covering 40 languages for generation, explanation, and completion tasks. FullStackBench~\citep{fullstackbench} assesses LLMs in realistic, multi-domain scenarios across 16 languages, employing a novel execution environment. Besides, some works propose more sophisticated evaluation frameworks. For example, ArtifactsBench~\citep{zhang2025artifactsbench} introduces automated multimodal evaluation for visual code generation, while CodeCriticBench~\citep{zhang2025codecriticbenchholisticcodecritique} focuses on holistic code critique evaluation. A common thread across these benchmarks is their reliance on labor-intensive manual curation for collecting problems, authoring ground-truth solutions, and designing test cases.

In contrast, our proposed work introduces a fully automated, dynamically constructed evaluation system. Unlike traditional benchmarks, it eliminates manual intervention, enabling a more scalable, consistent, and continuously evolving evaluation. Furthermore, it is designed to maintain a balanced difficulty distribution and support multiple languages, addressing the increasing demand for evaluating LLMs in diverse programming languages and contexts.


\subsection{Code Data Synthesis}
To reduce dependence on manually curated data, a growing body of research explores automatic data synthesis to augment the training of Code LLMs~\citep{luo2024wizardcoder,magicoder,zheng-etal-2024-opencodeinterpreter,wu2024inversecoderselfimprovinginstructiontunedcode,yu2024wavecoderwidespreadversatileenhancement,ahmad2025opencodeinstructlargescaleinstructiontuning,xu2025kodcode}. For instance, Evol-Instruct~\citep{luo2024wizardcoder} uses heuristic prompts to guide LLMs in evolving existing programming problems, thereby increasing their diversity and difficulty. OSS-Instruct~\citep{magicoder} prompts LLMs to generate new coding problems and solutions from raw, open-source code snippets. KodCode~\citep{xu2025kodcode} synthesizes a broad spectrum of Python coding tasks—including questions, solutions, and test cases—and ensures correctness through a systematic self-verification procedure. Some other methods focus on model self-improvement~\citep{wu2024inversecoder,NEURIPS2024selfcodealign,chen-etal-2025-revisit,zhou2025refinecoderiterativeimprovinglarge}. Inverse-Instruct~\citep{wu2024inversecoder} is a self-improvement technique that generates new instructions by "back-translating" code from an LLM's own training set, reducing the need to distill from more powerful proprietary models. Similarly, SelfCodeAlign~\citep{NEURIPS2024selfcodealign} introduces a pipeline for self-aligning code LLMs without extensive human annotation, using the same base model for both data generation and validation. Collectively, these data synthesis methods significantly reduce the reliance on manual curation and enable the continuous expansion of the problem space for training. Our work extends this paradigm of automation from data augmentation to the benchmark creation process. By leveraging extensive LLM-sandbox interaction, our pipeline not only automates the synthesis of verifiable test problems but can also be naturally repurposed for synthesizing high-quality training datasets. 

